% Chapter 2

\section{公式插图表格}

\subsection{公式的使用}
在文中引用公式可以这么写：$a^2+b^2=c^2$这是勾股定理，他还可以表示为$c=\sqrt{a^2+b^2}$，还可以让公式单独一段并且加上编号。注意，公式前请不要空行。
\begin{equation}
	\sin^2{\theta}+\cos^2{\theta}=1
	\label{eq:pingfanghe}
\end{equation}

还可以通过添加标签在正文中引用公式，如式\eqref{eq:pingfanghe}。

我们还可以轻松打出一个漂亮的矩阵：
\begin{equation}
	\mathbf{A}=
	\left[\begin{matrix}
			1  & 2  & 3  & 4  \\
			11 & 22 & 33 & 44 \\
		\end{matrix}\right] \times
	\left[\begin{matrix}
			22 & 24 \\
			32 & 34 \\
			42 & 44 \\
			52 & 54 \\
		\end{matrix}\right]
\end{equation}

或者多行对齐的公式：
\begin{equation}
	\begin{aligned}
		f_1(x) & =(x+y)^2     \\
		       & =x^2+2xy+y^2
	\end{aligned}
\end{equation}


\subsection{插图的使用}

\LaTeX 环境下可以使用常见的图片格式：JPEG、PNG、PDF、EPS等。当然也可以使用\LaTeX 直接绘制矢量图形，可以参考pgf/tikz等包中的相关内容。需要注意的是，无论采用什么方式绘制图形，首先考虑的是图片的清晰程度以及图片的可理解性，过于不清晰的图片将可能会浪费很多时间。

图示例如下：

\begin{figure}[!htb]
	\centering
	\includegraphics[width=0.3\textwidth]
	{figures/whulogo.png}\\
	\bicaption{插图示例}{Example of Figure}
	\label{fig:whu}
\end{figure}

\verb|[htbp]|选项分别是此处、页顶、页底、独立一页。\verb|[width=\textwidth]|让图片占满整行，或\verb|[width=2cm]|直接设置宽度。可以随时在文中进行引用，如图~\ref{fig:whu}，建议缩放时保持图像的宽高比不变。

\subsection{表格的使用}

表格的输入可能会比较麻烦，可以使用在线的工具，如~\href{https://www.tablesgenerator.com/}{Tables Generator}~能便捷的创建表格，也可以使用离线的工具，如~\href{https://ctan.org/pkg/excel2latex}{Excel2LaTeX}~支持从Excel表格转换成\LaTeX{}表格。\href{https://en.wikibooks.org/wiki/LaTeX/Tables}{LaTeX/Tables}~上及~\href{https://www.tug.org/pracjourn/2007-1/mori/mori.pdf}{Tables in LaTeX}~也有更多的示例能够参考。

\subsubsection{普通表格}
下面是一些普通表格的示例：

\begin{table}[ht]
	\centering
	\renewcommand\arraystretch{1.5}  % 扩大间距
	\bicaption{简单表格}{Simple Table}
	\label{tab:1}
	\begin{tabular}{|l|c|r|}
		\hline
		我是 & 一只 & 普通 \\
		\hline
		的   & 表格 & 呀   \\
		\hline
	\end{tabular}
\end{table}

\begin{table}[h!]
	\centering
	\bicaption{文献类型和标志代码}{Referrence type and identification code}
	\label{tab:2}
	\renewcommand\arraystretch{1.5}  % 扩大间距
	\newcolumntype{I}{!{\vrule width 2pt}}  % 加宽竖线
	\begin{tabular}{ccIcc}
		\toprule[2pt]
		{参考文献} & {文献类型} & {参考文献} & {文献类型} \\
		{类型}     & {标识}     & {类型}     & {标识}     \\
		\hline
		专著       & M          & 学位论文   & D          \\
		会议录     & C          & 报告       & R          \\
		期刊       & J          & 标准       & S          \\
		报纸       & N          & 专利       & P          \\
		汇编       & G          & 档案       & A          \\
		计算机程序 & CP         & 电子公告   & EB         \\
		数据库     & DB         & 美图       & CM         \\
		数据集     & DS         & 其他       & Z          \\
		\bottomrule[2pt]
	\end{tabular}
\end{table}

\begin{table}[ht]
	\centering
	\renewcommand\arraystretch{1.5}  % 扩大间距
	\bicaption{一般三线表}{Table with Three Rows}
	\label{tab:3}
	\begin{tabular}{ccc}
		\hline
		姓名 & 学号 & 性别 \\
		\hline
		张三 & 001  & 男   \\
		李四 & 002  & 女   \\
		\hline
	\end{tabular}
\end{table}

\subsubsection{统计表格 \& 跨栏表格和图片}
要创建占满整个文字宽度的表格需要使用到tabularx，如不需要，使用tabular就行。引用表格与其它引用一样，只需要：表~\ref{tab:4}，统计表格一般是三线表形式。

跨栏表格和图片请使用\verb|table*|和\verb|figure*|。

% 跨栏表格宽度占整页，为单栏样式，使用由cuted提供的strip代替table（跨栏图片用strip代替figure），\verb|\caption{标题}|改为\verb|\captionof{table}{标题}|（图片则是\verb|\captionof{figure}{标题}|）。

\begin{table*}
	\centering
	\bicaption{统计数据表格}{Table of Statistic Data}
	\label{tab:4}
	\begin{tabularx}{\textwidth}{CCCC}
		\toprule
		序号 & 年龄 & 身高   & 体重  \\
		\midrule
		1    & 14   & 156    & 42    \\
		2    & 16   & 158    & 45    \\
		3    & 14   & 162    & 48    \\
		4    & 15   & 163    & 50    \\
		\cmidrule{2-4} %添加2-4列的中线
		平均 & 15   & 159.75 & 46.25 \\
		\bottomrule
	\end{tabularx}
\end{table*}

\subsection{列表的使用}
下面演示了创建有序及无序列表，如需其它样式，\href{https://www.latex-tutorial.com/tutorials/lists/}{LaTeX Lists}~上有更多的示例。

\subsubsection{有序列表}
这是一个计数的列表
\begin{enumerate}
	\item 第一项
	      \begin{enumerate}
		      \item 第一项中的第一项
		      \item 第一项中的第二项
	      \end{enumerate}
	\item 第二项
	      \begin{enumerate}[label=(\roman*)]
		      \item 第一项中的第一项
		      \item 第一项中的第二项
	      \end{enumerate}
	\item 第三项
\end{enumerate}

\subsubsection{不计数列表}
这是一个不计数的列表
\begin{itemize}
	\item 第一项
	      \begin{itemize}
		      \item 第一项中的第一项
		      \item 第一项中的第二项
	      \end{itemize}
	\item 第二项
	\item 第三项
\end{itemize}

\subsection{定理的使用}
\begin{theorem}
	设向量$\boldsymbol a\neq\boldsymbol 0$，那么向量$\boldsymbol b//\boldsymbol a$的充分必要条件是：存在唯一的实数$\lambda$，使$\boldsymbol b=\lambda \boldsymbol a$。
\end{theorem}
\begin{definition}
	这是一条定义。
\end{definition}
\begin{lemma}
	这是一条引理。
\end{lemma}
\begin{corollary}
	对数轴上任意一点$P$，轴上有向线段$\vec {OP}$都可唯一地表示为点$P$的坐标与轴上单位向量$\boldsymbol e_u$的乘积：$\vec {OP}=u \boldsymbol e_u$。
\end{corollary}
\begin{proposition}
	这是一条性质。
\end{proposition}
\begin{example}
	这是一条例。
\end{example}
\begin{remark}
	这是一条注。
\end{remark}